% ============================================================
%  Macros pròpies de la teoria de categories, locals al volum.
%
%  El fitxer compartit ../common/macros.tex barreja notació
%  genèrica (\R, \N, \mcd, decimals periòdics...) amb la de
%  categories. Aquest fitxer recull, en un sol lloc i documentada,
%  només la notació que aquest llibre fa servir, de manera que el
%  volum no depèn de l'organització del fitxer compartit.
%
%  S'usa \providecommand perquè, si macros.tex ja ha definit una
%  macro, no es produeixi cap col·lisió: la definició compartida i
%  la local són idèntiques. Si algun dia es depura macros.tex
%  traient-ne la part de categories, aquest fitxer la garanteix.
% ============================================================

\providecommand{\cat}[1]{\mathbf{#1}}          % noms de categories: \cat{Set}
\providecommand{\id}{\operatorname{id}}         % identitat
\providecommand{\Hom}{\operatorname{Hom}}       % conjunts de morfismes
\providecommand{\op}{^{\mathrm{op}}}            % categoria oposada
\providecommand{\Obj}{\operatorname{Obj}}       % col·lecció d'objectes
\providecommand{\comp}{\circ}                    % composició
\providecommand{\yon}{\mathrm{y}}                % plançó de Yoneda

\providecommand{\dom}{\operatorname{dom}}
\providecommand{\cod}{\operatorname{cod}}

% Predicats relacionals de l'apèndix d'axiomàtica relacional (Dom, Cod, Com,
% Id). S'escriuen amb versaletes matemàtiques per distingir-los clarament dels
% operadors funcionals \dom, \cod: aquí són relacions primitives, no funcions.
\providecommand{\Dom}{\mathsf{Dom}}
\providecommand{\Cod}{\mathsf{Cod}}
\providecommand{\Com}{\mathsf{Com}}
\providecommand{\Id}{\mathsf{Id}}
\providecommand{\Free}{\operatorname{Free}}
\providecommand{\Nat}{\operatorname{Nat}}        % transformacions naturals com a homconjunt
\providecommand{\Sub}{\operatorname{Sub}}        % preordre de subobjectes
\providecommand{\ev}{\operatorname{ev}}          % morfisme d'avaluació (exponencials)
\providecommand{\Imatge}{\operatorname{Im}}      % imatge d'un morfisme

% Claudàtors semàntics [[-]] sense dependre de stmaryrd:
\providecommand{\llbracket}{\mathopen{[\![}}
\providecommand{\rrbracket}{\mathclose{]\!]}}

% colimit amb subíndex a sota (com \lim), només si encara no existeix:
\makeatletter
\@ifundefined{colimop}{\DeclareMathOperator*{\colimop}{colim}}{}
\makeatother